% Part: first-order-logic
% Chapter: sequent-calculus
% Section: proving-things

\documentclass[../../../include/open-logic-section]{subfiles}

\begin{document}

\iftag{FOL}
      {\olfileid{fol}{seq}{pro}}
      {\olfileid{pl}{seq}{pro}}

\olsection{\usetoken{P}{derivation}示例}

\begin{ex}
为矢列 $!A \land !B \Sequent !A$ 构造一个 $\Log{LK}$-推导。

我们先把这个想要得到的末端矢列写在推导的底部。
\begin{prooftree}
\AxiomC{}
\UnaryInf$!A\land !B \fCenter !A$
\end{prooftree}
接下来，我们需要弄清楚哪种推理能得到这种形式的下端矢列。它可能是结构规则，但一个好的思路是从寻找逻辑规则开始。末端矢列中出现的唯一逻辑算子是 $\land$，所以我们找的是 $\land$ 规则；又因为 $\land$ 符号出现在前件中，我们考虑的是 \LeftR{\land} 规则。
\begin{prooftree}
\AxiomC{}
\RightLabel{\LeftR{\land}}
\UnaryInf$!A\land !B \fCenter !A$
\end{prooftree}
对于这个 \LeftR{\land} 推理的上矢列，有两种可能的选择：上矢列可以是 $!A \Sequent !A$，也可以是 $!B \Sequent !A$。显然，$!A \Sequent !A$ 是一个初始矢列（这是好事），而 $!B \Sequent !A$ 通常不可推导。我们填入上矢列：
\begin{prooftree}
\Axiom$!A \fCenter !A$
\RightLabel{\LeftR{\land}}
\UnaryInf$!A\land !B \fCenter !A$
\end{prooftree}
现在我们得到了矢列 $!A \land !B \Sequent !A$ 的一个正确的 $\Log{LK}$-推导。
\end{ex}

\begin{ex}
为矢列 $\lnot !A \lor !B
\Sequent !A \lif !B$ 构造一个 $\Log{LK}$-推导。

首先将想要的末端矢列写在推导的底部。
\begin{prooftree}
\AxiomC{}
\UnaryInf$\lnot !A \lor !B \fCenter !A \lif !B$
\end{prooftree}
为了找到能得到该末端矢列的逻辑规则，我们先看末端矢列中的逻辑联结词：$\lnot$、$\lor$ 和 $\lif$。我们此刻只关心 $\lor$ 和 $\lif$，因为它们是末端矢列中语句的主联结词，而 $\lnot$ 在另一个联结词的辖域内，所以我们稍后再处理它。因此，最后一步推理可用的逻辑规则是 \LeftR{\lor} 规则和 \RightR{\lif} 规则。实际上我们可以选其中任何一个，但我们选择 \RightR{\lif} 规则（如果非要找个理由的话，那就是它允许我们推迟拆分成两个分支）。根据能得到下端矢列的 \RightR{\lif} 推理形式，它必须形如：
\begin{prooftree}
\AxiomC{}
\UnaryInf$ !A, \lnot !A \lor !B \fCenter !B $
\RightLabel{\RightR{\lif}} \UnaryInf$ \lnot !A \lor !B \fCenter !A \lif !B $
\end{prooftree}

如果我们将 $\lnot !A \lor !B$ 移到前件的最前面，就可以应用 \LeftR{\lor} 规则。根据该规则的图式，这里必须拆分成如下两个上端矢列：
\begin{prooftree}
\AxiomC{}
\UnaryInf$\lnot !A, !A \fCenter !B$
\AxiomC{}
\UnaryInf$!B, !A \fCenter !B$
\RightLabel{\LeftR{\lor}}
\BinaryInf$ \lnot !A \lor !B, !A \fCenter !B $
\RightLabel{\RightR{\Exchange}}
\UnaryInf$ !A, \lnot !A \lor !B \fCenter !B $
\RightLabel{\RightR{\lif}}
\UnaryInf$ \lnot !A \lor !B \fCenter !A \lif !B $
\end{prooftree}
记住，我们的目标是向上回溯到初始矢列；我们似乎已经非常接近了！右分支只需一次弱化和一次交换就能变成初始矢列，然后该分支就完成了：
\begin{prooftree}
\AxiomC{}
\UnaryInf$\lnot !A, !A \fCenter !B$
\Axiom$!B \fCenter !B$
\RightLabel{\LeftR{\Weakening}}
\UnaryInf$!A, !B \fCenter !B$
\RightLabel{\LeftR{\Exchange}}
\UnaryInf$!B, !A \fCenter !B$
\RightLabel{\LeftR{\lor}}
\BinaryInf$\lnot !A \lor !B, !A \fCenter !B $
\RightLabel{\RightR{\Exchange}}
\UnaryInf$ !A, \lnot !A \lor !B \fCenter !B $
\RightLabel{\RightR{\lif}}
\UnaryInf$ \lnot !A \lor !B \fCenter !A \lif !B $
\end{prooftree}

现在来看左分支。其中各语句唯一的逻辑联结词是前件中的 $\lnot$ 符号，因此我们需要应用 \LeftR{\lnot} 规则的一个实例。
\begin{prooftree}
\AxiomC{}
\UnaryInf$ !A \fCenter !B, !A$
\RightLabel{\LeftR{\lnot}}
\UnaryInf$\lnot !A, !A \fCenter !B$
\Axiom$!B \fCenter !B$
\RightLabel{\LeftR{\Weakening}}
\UnaryInf$!A, !B \fCenter !B$
\RightLabel{\LeftR{\Exchange}}
\UnaryInf$!B, !A \fCenter !B$
\RightLabel{\LeftR{\lor}}
\BinaryInf$\lnot !A \lor !B, !A \fCenter !B $
\RightLabel{\RightR{\Exchange}}
\UnaryInf$ !A, \lnot !A \lor !B \fCenter !B $
\RightLabel{\RightR{\lif}}
\UnaryInf$ \lnot !A \lor !B \fCenter !A \lif !B $
\end{prooftree}
类似于完成右分支的方式，左分支也只需一次弱化和一次交换即可完成。
\begin{prooftree}
\Axiom$!A \fCenter !A$
\RightLabel{\RightR{\Weakening}}
\UnaryInf$ !A \fCenter !A, !B$
\RightLabel{\RightR{\Exchange}}
\UnaryInf$ !A \fCenter !B, !A$
\RightLabel{\LeftR{\lnot}}
\UnaryInf$\lnot !A, !A \fCenter !B$
\Axiom$!B \fCenter !B$
\RightLabel{\LeftR{\Weakening}}
\UnaryInf$!A, !B \fCenter !B$
\RightLabel{\LeftR{\Exchange}}
\UnaryInf$!B, !A \fCenter !B$
\RightLabel{\LeftR{\lor}}
\BinaryInf$\lnot !A \lor !B, !A \fCenter !B $
\RightLabel{\RightR{\Exchange}}
\UnaryInf$ !A, \lnot !A \lor !B \fCenter !B $
\RightLabel{\RightR{\lif}}
\UnaryInf$ \lnot !A \lor !B \fCenter !A \lif !B $
\end{prooftree}
\end{ex}

\begin{ex}
为矢列 $\lnot !A \lor \lnot !B
\Sequent \lnot (!A \land !B)$ 构造一个 $\Log{LK}$-推导。

运用上面的技巧，我们首先将想要的末端矢列写在底部。
\begin{prooftree}
\AxiomC{}
\UnaryInf$ \lnot !A \lor \lnot !B \fCenter \lnot (!A \land !B) $
\end{prooftree}
末端矢列中各语句可用的主联结词是 $\lor$ 符号和 $\lnot$ 符号。此时应用 \LeftR{\lor} 规则或 \RightR{\lnot} 规则都可以，但我们从 \RightR{\lnot} 规则开始，因为它能暂时避免分裂成两个分支：
\begin{prooftree}
\AxiomC{}
\UnaryInf$!A \land !B, \lnot !A \lor \lnot !B \fCenter $
\RightLabel{\RightR{\lnot}}
\UnaryInf$\lnot !A \lor \lnot !B \fCenter \lnot (!A \land !B)$
\end{prooftree}
现在我们面临一个选择：考虑 \LeftR{\land} 规则还是 \LeftR{\lor} 规则。让我们看看如果应用 \LeftR{\land} 规则会发生什么：我们可以选择从矢列 $!A,
\lnot !A \lor !B \Sequent \quad$ 或矢列 $!B, \lnot !A
\lor !B \Sequent \quad$ 开始。由于推导关于 $!A$ 和 $!B$ 是对称的，我们就选前者：
\begin{prooftree}
\AxiomC{}
\UnaryInf$!A, \lnot !A \lor \lnot !B \fCenter $
\RightLabel{\LeftR{\land}}
\UnaryInf$!A \land !B, \lnot !A \lor \lnot !B \fCenter $
\RightLabel{\RightR{\lnot}}
\UnaryInf$\lnot !A \lor \lnot !B \fCenter \lnot (!A \land !B)$
\end{prooftree}
继续填入推导过程，我们发现遇到了问题：
\begin{prooftree}
\Axiom$!A \fCenter !A$
\RightLabel{\LeftR{\lnot}}
\UnaryInf$ \lnot !A, !A \fCenter$
\AxiomC{}
\RightLabel{?}
\UnaryInf$!A \fCenter !B$
\RightLabel{\LeftR{\lnot}}
\UnaryInf$ \lnot !B, !A \fCenter$
\RightLabel{\LeftR{\lor}}
\BinaryInf$\lnot !A \lor \lnot !B, !A \fCenter $
\RightLabel{\LeftR{\Exchange}}
\UnaryInf$!A, \lnot !A \lor \lnot !B \fCenter $
\RightLabel{\LeftR{\land}}
\UnaryInf$!A \land !B, \lnot !A \lor \lnot !B \fCenter $
\RightLabel{\RightR{\lnot}}
\UnaryInf$\lnot !A \lor \lnot !B \fCenter \lnot (!A \land !B)$
\end{prooftree}
右分支的顶端无法再进一步归约，也无法通过结构推理变成初始矢列，所以这条路径是不对的。因此很明显，刚才应用 \LeftR{\land} 规则是个错误。回到之前的状态，转而执行 \LeftR{\lor} 规则，我们得到：
\begin{prooftree}
\AxiomC{}
\UnaryInf$\lnot !A, !A \land !B \fCenter $

\AxiomC{}
\UnaryInf$\lnot !B, !A \land !B \fCenter $

\RightLabel{\LeftR{\lor}}
\BinaryInf$\lnot !A \lor \lnot !B, !A \land !B \fCenter $
\RightLabel{\LeftR{\Exchange}}
\UnaryInf$!A \land !B, \lnot !A \lor \lnot !B \fCenter $
\RightLabel{\RightR{\lnot}}
\UnaryInf$\lnot !A \lor \lnot !B \fCenter \lnot (!A \land !B)$
\end{prooftree}
像之前那样完成每个分支，我们得到：
\begin{prooftree}
\Axiom$ !A \fCenter!A$
\RightLabel{\LeftR{\land}}
\UnaryInf$!A \land !B \fCenter !A$
\RightLabel{\LeftR{\lnot}}
\UnaryInf$\lnot !A, !A \land !B \fCenter $

\Axiom$ !B \fCenter !B$
\RightLabel{\LeftR{\land}}
\UnaryInf$!A \land !B \fCenter !B$
\RightLabel{\LeftR{\lnot}}
\UnaryInf$\lnot !B, !A \land !B \fCenter $

\RightLabel{\LeftR{\lor}}
\BinaryInf$\lnot !A \lor \lnot !B, !A \land !B \fCenter $
\RightLabel{\LeftR{\Exchange}}
\UnaryInf$!A \land !B, \lnot !A \lor \lnot !B \fCenter $
\RightLabel{\RightR{\lnot}}
\UnaryInf$\lnot !A \lor \lnot !B \fCenter \lnot (!A \land !B)$
\end{prooftree}
（在这些步骤中，我们本可以将 $\land$ 规则排在 $\lnot$ 规则的下面执行，同样能得到正确的推导）。
\end{ex}

\begin{ex}
到目前为止我们还没有使用过收缩规则，但它有时是必需的。这里就有一个例子。假设我们想证明 $\quad \Sequent !A \lor \lnot !A$。逆向应用 $\RightR{\lor}$ 规则会得到以下两种推导之一：
\begin{prooftree}
\AxiomC{}
\UnaryInf$ \fCenter !A$
\RightLabel{\RightR{\lor}}
\UnaryInf$ \fCenter !A \lor \lnot !A$
\DisplayProof\qquad\bottomAlignProof
\AxiomC{}
\UnaryInf$!A \fCenter $
\RightLabel{\RightR{\lnot}}
\UnaryInf$ \fCenter \lnot !A$
\RightLabel{\RightR{\lor}}
\UnaryInf$ \fCenter !A \lor \lnot !A$
\end{prooftree}
当然，这两个推导的顶端都不是初始矢列。这里的技巧在于认识到收缩规则允许我们将一个语句的两份副本合并为一份——当我们在搜索证明（即从下往上进行）时，可以在前提中保留一份 $!A \lor \lnot !A$ 的副本，例如：
\begin{prooftree}
\AxiomC{}
\UnaryInf$ \fCenter !A \lor \lnot !A, !A$
\RightLabel{\RightR{\lor}}
\UnaryInf$ \fCenter !A \lor \lnot !A, !A \lor \lnot !A$
\RightLabel{\RightR{\Contraction}}
\UnaryInf$ \fCenter !A \lor \lnot !A$
\end{prooftree}
现在我们可以第二次应用 $\RightR{\lor}$ 规则，同时得到~$\lnot !A$，这就引向了一个完整的推导。
\begin{prooftree}
\Axiom$!A \fCenter !A$
\RightLabel{\RightR{\lnot}}
\UnaryInf$\fCenter !A, \lnot !A$
\RightLabel{\RightR{\lor}}
\UnaryInf$\fCenter !A, !A \lor \lnot !A$
\RightLabel{\RightR{\Exchange}}
\UnaryInf$ \fCenter !A \lor \lnot !A, !A$
\RightLabel{\RightR{\lor}}
\UnaryInf$ \fCenter !A \lor \lnot !A, !A \lor \lnot !A$
\RightLabel{\RightR{\Contraction}}
\UnaryInf$ \fCenter !A \lor \lnot !A$
\end{prooftree}
\end{ex}

\begin{prob}
给出以下矢列的推导：
\begin{enumerate}
\item $!A \land (!B \land !C) \Sequent (!A \land !B) \land !C$。
\item $!A \lor (!B \lor !C) \Sequent (!A \lor !B) \lor !C$。
\item $!A \lif (!B \lif !C) \Sequent !B \lif (!A \lif !C)$。
\item $!A \Sequent \lnot\lnot !A$。
\end{enumerate}
\end{prob}

\begin{prob}
给出下列矢列的推导：
\begin{enumerate}
\item $(!A \lor !B) \lif !C \Sequent !A \lif !C$.
\item $(!A \lif !C) \land (!B \lif !C) \Sequent (!A \lor !B) \lif !C$.
\item $\Sequent \lnot(!A \land \lnot !A)$.
\item $!B \lif !A \Sequent \lnot !A \lif \lnot !B$.
\item $\Sequent (!A \lif \lnot !A) \lif \lnot !A$.
\item $\Sequent \lnot(!A \lif !B) \lif \lnot !B$.
\item $!A \lif !C \Sequent \lnot (!A \land \lnot !C)$.
\item $!A \land \lnot !C \Sequent \lnot (!A \lif !C)$.
\item $!A \lor !B, \lnot !B \Sequent !A$.
\item $\lnot !A \lor \lnot !B \Sequent \lnot(!A \land !B)$.
\item $\Sequent (\lnot !A \land \lnot !B) \lif\lnot(!A \lor !B)$.
\item $\Sequent \lnot(!A \lor !B) \lif (\lnot !A \land \lnot !B)$.
\end{enumerate}
\end{prob}

\begin{prob}
给出下列矢列的推导：
\begin{enumerate}
\item $\lnot(!A \lif !B) \Sequent !A$.
\item $\lnot(!A \land !B) \Sequent \lnot !A \lor \lnot !B$.
\item $!A \lif !B \Sequent \lnot !A \lor !B$.
\item $\Sequent \lnot \lnot !A \lif !A$.
\item $!A \lif !B, \lnot !A \lif !B \Sequent !B$.
\item $(!A \land !B) \lif !C \Sequent (!A \lif !C) \lor (!B \lif !C)$.
\item $(!A \lif !B) \lif !A \Sequent !A$.
\item $\Sequent (!A \lif !B) \lor (!B \lif !C)$.
\end{enumerate}
（这些推导都需要用到 $\RightR{\Contraction}$~规则。）
\end{prob}


\end{document}